% ============================================================
%  演讲稿模板 —— 通用学术/组会讲稿（LaTeX article）
%  编译方式: XeLaTeX 或 pdfLaTeX 均可
%  用法: 替换「【】」标记的内容，删除示例段落
% ============================================================
\documentclass[12pt,a4paper]{article}

% ---------- 页面与中文支持 ----------
\usepackage[margin=2cm]{geometry}
\usepackage{ctex}                % 中文支持（XeLaTeX 下推荐）
\usepackage{amsmath,amssymb}
\usepackage{enumitem}
\usepackage{graphicx}
\usepackage{hyperref}
\usepackage{xcolor}

\hypersetup{colorlinks=true, linkcolor=blue, urlcolor=blue}

% ============================================================
%	标题区 —— 请填写
% ============================================================
\title{\textbf{【汇报/演讲标题】}}
\author{【姓名】 \\ \small 【单位/学院】}
\date{\today}

\begin{document}
\maketitle
\vspace{-1.2cm}

% ============================================================
%	开场 —— 问候 + 自我介绍 + 点题
% ============================================================
\section*{开场（30 秒）}
【各位老师/同学，大家好！我是……，今天我汇报的题目是……。】

本报告将从以下三个部分展开：第一部分介绍【背景与问题】；第二部分阐述【方法与原理】；第三部分给出【结果与结论】。

% ============================================================
%	一、背景与问题
% ============================================================
\section{背景与问题}
\subsection{为什么研究这个问题}
【研究动机：现有方法存在哪些不足？本工作试图解决什么问题？】
\begin{itemize}
    \item 问题一：……
    \item 问题二：……
\end{itemize}

\subsection{相关工作概述}
【简要梳理已有研究，突出本文差异化。】
\begin{itemize}
    \item 方法 A（作者，年份）：……优点/缺点
    \item 方法 B（作者，年份）：……优点/缺点
\end{itemize}

% ============================================================
%	二、方法与原理（核心部分，讲解最详细）
% ============================================================
\section{方法}
\subsection{总体思路}
【一两句话概括方法的核心思想。】

\subsection{数学原理}
【核心公式用 equation 环境，逐个解释每个符号的含义。】
\begin{equation}
    y = \mathbf{W}^{\top} \mathbf{x} + b
    \label{eq:core}
\end{equation}
其中 $\mathbf{x} \in \mathbb{R}^n$ 为输入特征，$\mathbf{W}$ 为权重矩阵，$b$ 为偏置。式 \eqref{eq:core} 表示……

\subsection{算法流程}
\begin{enumerate}[leftmargin=2em]
    \item 步骤一：……
    \item 步骤二：……
    \item 步骤三：……
\end{enumerate}

\subsection{可视化/图表}
【若用图说明，将图片放入同目录并取消注释。】
% \begin{figure}[htbp]
%     \centering
%     \includegraphics[width=0.6\textwidth]{figure.png}
%     \caption{图注：方法示意图/实验结果}
%     \label{fig:method}
% \end{figure}

% ============================================================
%	三、结果与分析
% ============================================================
\section{结果与分析}
\subsection{实验设置}
【数据集、评价指标、对比方法、超参数。】
\begin{table}[htbp]
    \centering
    \begin{tabular}{lccc}
        \hline
        \textbf{方法} & \textbf{指标 A} & \textbf{指标 B} & \textbf{指标 C} \\
        \hline
        基线方法 & 数值 & 数值 & 数值 \\
        本文方法 & 数值 & 数值 & 数值 \\
        \hline
    \end{tabular}
    \caption{实验结果对比表}
\end{table}

\subsection{主要结论}
\begin{itemize}
    \item 结论一：……
    \item 结论二：……
\end{itemize}

% ============================================================
%	收尾 —— 总结 + 致谢
% ============================================================
\section*{总结与展望}
【用 2-3 句话总结贡献，再简要说明未来工作方向。】

\vspace{0.5cm}
\textbf{我的汇报到此结束，感谢各位老师、同学的聆听，恳请批评指正！}

\end{document}
